% Section 8.4, from lectures/L08/README.md: what you should be able to do, the questions to test
% yourself with, and the next lecture. The closing paragraph is from the README's reference list.

\needspace{10\baselineskip}
\section{Sammanfattning}
\label{c8:sec:summary}

\needspace{6\baselineskip}
\subsection*{Det här ska du kunna nu}
Efter det här kapitlet ska du kunna:
\begin{itemize}
  \item Addera och subtrahera binärt för hand, och säga vad ett register och \code{C} innehåller
    efteråt.
  \item Förklara aritmetisk rundgång, och varför \code{inc} och \code{dec} inte ändrar \code{C}.
  \item Omvandla mellan ett negativt tal och dess tvåkomplement, och läsa samma byte med och utan
    tecken.
  \item Förutsäga \code{C}, \code{Z}, \code{N} och \code{V} efter en addition eller subtraktion,
    och förklara skillnaden mellan \code{C} och \code{V}.
  \item Välja mask och instruktion för att ettställa, nollställa, invertera eller testa enskilda
    bitar.
  \item Använda skift för att multiplicera och dividera med 2, och förklara skillnaden mellan
    \code{lsr} och \code{asr}.
  \item Göra en port till utport, och tända och släcka enskilda stift med \code{out}, \code{sbi}
    och \code{cbi}.
\end{itemize}

\needspace{6\baselineskip}
\subsection*{Frågor att testa dig själv med}
\begin{enumerate}
  \item \code{r16} är 255 och \code{C} är 0. Vad är \code{r16} och \code{C} efter \code{inc r16}?
    Och efter \code{add r16, r17} med \code{r17}~=~1?
  \item Varför finns det ingen \code{addi}, och hur adderar man 5 till \code{r16} ändå?
  \item Byten \code{0xF6}: vilket tal är den utan tecken, och vilket med tecken?
  \item \code{0x80 - 0x01} ger \code{C = 0} men \code{V = 1}. Vad säger det om svaret, utan tecken
    och med tecken?
  \item Vilken instruktion och mask nollställer bit~7 i \code{r16} utan att röra de andra bitarna?
  \item Varför är ett av nio steg mörkt i det rinnande ljuset med \code{rol}?
  \item \code{out PORTB, r16} och \code{sbi PORTB, 3}: vad skiljer dem åt, när det gäller de andra
    sju stiften?
\end{enumerate}

Referensbladet i bilaga~\ref{app:avr} visar vilka flaggor varje instruktion ändrar.

\needspace{6\baselineskip}
\subsection*{Nästa kapitel}
Kapitel~\ref{c9:ch} handlar om programflöde, loopar och tidsfördröjningar:
\begin{itemize}
  \item Flödesplaner: att rita ett program innan det skrivs.
  \item Jämförelser och villkorliga hopp: det flaggorna är till för.
  \item Loopar som går ett bestämt antal varv.
  \item Tidsfördröjningar, beräknade exakt från klockcykler och kontrollerade med simulatorns
    stoppur.
\end{itemize}
